\section{Conclusion and Further Work}

The hidden Hodge operator determines the recovery target. An explicitly
presented harmonic lattice gives a precise discrete-CVP comparison, while
the constructed common-input examples prevent uniform exact recovery from
the specified public data. Exact-component observations and full projective
observations expose different information, as the two observation results
make precise.

\subsection{Further Questions}

The comparison suggests studying encoded HCP families with explicit
distributions on hidden presentations and Hodge-structured noise. Their
computational difficulty must be assessed within that specified model.

Discrete Morse reduction for cellular sheaf cohomology
\cite{curryGhristNanda2015morse} suggests a computational refinement.
A chain-homotopy equivalence alone does not identify the original harmonic
representative or preserve its distance. An application to the original
recovery target needs explicit retained lifts and a metric comparison.

The constructed classifier proves the named homological-test fragment in the
real cochain model. Extending it to the full first-order fibrational semantics,
constructing richer families with a declared topological realization, and
analyzing transported public observations are separate mathematical tasks.
The selected results remain the exact Hodge/CVP comparisons and observation
boundaries of this manuscript.
